% =====================================================================
\chapter*{Preface}
\addcontentsline{toc}{chapter}{Preface}
\markboth{Preface}{Preface}
% =====================================================================

\lettrine[lines=2,findent=2pt,nindent=0pt]{\color{inkred}E}{very} model we train, from a small linear regression to a large language model, learns in the same way: an optimizer follows the gradient of a loss, one step at a time, hoping to reach its minimum. This book is about those steps.

It started, like the first \emph{Enchiridion}, as handwritten notes. From May to June 2026, I wrote by hand the theory of the optimizers used in machine learning and deep learning: first the mathematical tools (convexity and L-smoothness), then Gradient Descent and its proof of convergence, and then every idea that was added to it to make it faster, more stable, or more robust: Stochastic Gradient Descent, Momentum, Nesterov, RMSProp, Adam, AdamW, and the other versions of Adam, until the learning rate schedules that decide how fast we move.

An \emph{enchiridion} is, in Greek, ``something that is held in the hand'': a handbook. This is the second one of the series, after \emph{The Enchiridion for Tree Based Algorithms Theory}.

\section*{Who made this book}

\begin{tcolorbox}[colback=paper,colframe=goldleaf,boxrule=0.8pt,arc=3pt,left=12pt,right=12pt]
\textbf{Author.} The content of this book (the explanations, the proofs, the examples and the way of telling them) was created by \textbf{Reda Louriki}, in handwritten notes.

\medskip
\textbf{Typesetting.} The handwritten notes were transcribed and turned into \LaTeX\ by an AI, \textbf{Claude Opus 5.5} (Anthropic). The author's wording and examples were kept; only grammatical errors and mathematical errors were corrected. The part of the last chapter that the notes left unfinished (the warm up) was completed in the same spirit.
\end{tcolorbox}

Every time a correction changes the meaning of what was written in the notes (a formula, an index, a number, a name, or a conclusion), it is signalled with a footnote marked \textsf{\bfseries\color{inkgold}Editor's note}, so you can always see what the notes said and why it was changed. The list of all these changes is gathered at the end of the book, in \emph{Notes on this Edition}.

\section*{How to read this book}

The notes were written with several colours of ink, and each colour had a job. The book keeps the same code:

\begin{center}
\renewcommand{\arraystretch}{1.3}
\begin{tabular}{l l}
{\color{inkred}\itshape Red} & the titles of the chapters and the warnings \warn\\
{\color{inkgreen}\itshape Green} & definitions, theorems, ``Problem'' and ``Solution''\\
{\color{inkgold}\itshape ``Gold''} & the motto of each optimizer, in one sentence\\
{\color{inkblue}\itshape Blue} & demonstrations and remarks\\
{\color{inkpurple}\itshape Purple} & the questions you should be asking yourself at that moment\\
\tikz[baseline=-0.6ex]\node[fill=goldpale!45,draw=goldleaf,inner sep=2pt,font=\itshape\color{inkgold}]{Gold boxes}; & the sentences worth remembering\\
\end{tabular}
\end{center}

The chapters follow the logic of the ideas: each optimizer is born from a problem of the previous one. Every chapter starts with that problem, and ends with its solution. So the best way to read this book is in order.

The prerequisites are some linear algebra (inner products, norms, eigenvalues), derivatives and gradients, some probability (expectation, variance), and the curiosity to ask ``why?'' at every line. The cover shows the idea of the whole book: the real trajectories of six of these optimizers, computed on the same loss landscape, from the same starting point, racing to the same minimum.

\bigskip
\begin{flushright}
\itshape Reda Louriki\\
2026
\end{flushright}
